\section*{What the learned body is for, and how to reward it}
The executor owns a build; the learned body acts only when the agent is hurt. A weighting of event kinds ($r=\phi^\top
w$) cannot express ``build this design'', so the body is not where building is learned. Its role is the short skills
the executor and the language model call: go to a place, gather an item, fight, flee. Three results bear on it: one
network can serve many reward weightings (E1); skills learned once combine well on new tasks (E2); and progress the
game can verify can be rewarded at every step without changing what is optimal, which is what makes a sparse goal
learnable (S). E1 and E2 run on Garnet MDPs $G(30,4,5)$ [Archibald et al.\ 1995] with $K=4$ reward kinds, $\gamma=0.9$, 20
seeds; S on a small build task, 20 seeds. Each experiment isolates one claim [Osband et al.\ 2020]; intervals are
percentile bootstraps of the interquartile mean over seeds [Agarwal et al.\ 2021]. Code: \texttt{research/bench.py}
(24\,s on one core).
